\section{Introduction}\label{sec:introduction}

A scaling limit for the metric and area of a random surface does not by
itself determine the motion of a particle on that surface. The electrical
conductances of the approximating graphs can change the limiting
diffusion, even when their measured metric spaces converge. A further
question is whether the time needed to see the continuum motion has a
single deterministic normalization.

We resolve this question for planar maps weighted by their numbers of
spanning trees. The geometric input is the metric-measure limit and the
contour reconstruction theory of \cite{MetricMaps}. We prove the additional
convergence of the conditional walk laws, with an exact linear clock.
The proof separates the identification of the electrical energy from the
control of the speed measure. This separation is useful because neither
part follows from metric-measure convergence alone.

\subsection{The model and the continuum process}

A rooted planar map is a finite connected graph embedded in the oriented
sphere, up to orientation-preserving homeomorphism, with a distinguished
oriented edge. Loops and multiple edges are allowed. For $n\ge1$, let
$(M_n,T_n)$ be uniform among rooted maps with $n$ edges decorated by one
of their spanning trees. Thus the marginal weight of $M_n$ is proportional
to its number of spanning trees; conditional on $M_n$, the tree $T_n$ is
uniform. Write $V_n$ for the vertex set and $d_n$ for graph distance using
all map edges, each of length one. Set
\begin{equation}\label{eq:degree-measure}
 \mu_n=\sum_{v\in V_n}\frac{\deg(v)}{2n}\,\delta_v,
\end{equation}
where each loop contributes two to the degree.

Conditional on the map, let $X^n$ start from $\mu_n$. At the times of an
independent rate-one Poisson process it chooses an incident half-edge
uniformly and moves to its other endpoint. Choosing a loop makes no
displacement. Its generator is
\begin{equation}\label{eq:walk-generator}
 (L_nf)(v)=\frac1{\deg(v)}
   \sum_{e\text{ a half-edge at }v}
       \bigl(f(\operatorname{other}(e))-f(v)\bigr).
\end{equation}
The measure $\mu_n$ is stationary and reversible for this walk.

Let $(S,h)$ be an ordinary unit-area $\sqrt2$-Liouville quantum sphere,
with its distinguished points forgotten. We use the ordinary fixed-area
sphere law of Duplantier, Miller, and Sheffield \cite{DMS}. Its quantum
area measure $\mu_h$ has total mass one. Write $D_h$ for its Liouville
quantum gravity metric in a fixed deterministic normalization. Identify
the underlying conformal sphere with the round sphere. The continuum
process $B^h$ is the conservative $\mu_h$-symmetric diffusion associated
with the closure on $L^2(\mu_h)$ of
\begin{equation}\label{eq:continuum-form}
 \cE_h(f,g)=\frac12\int_{S^2}
      \langle\nabla f,\nabla g\rangle_{\mathrm{round}}
                    \,\dd\vol_{\mathrm{round}},
 \qquad f,g\in C^\infty(S^2).
\end{equation}
The closed form uses quasi-continuous representatives. This is Liouville
Brownian motion, started with conditional law $\mu_h$; see
\cite{GRV,Berestycki} and the Dirichlet-form framework in \cite{FOT}.
Conformal invariance of the two-dimensional Dirichlet integral makes
\eqref{eq:continuum-form} independent of the conformal coordinate. Its
factor $1/2$ fixes the continuum clock.

For precision we fix the spatial normalization used throughout. For a
finite real random variable $Y$, let $m(Y)$ be the unique real solution
of $\E[\arctan(Y-m(Y))]=0$. Put
\begin{equation}\label{eq:spatial-normalization}
 \Delta_n=\max\{1,\diam(V_n,d_n)\},\qquad
 \Delta=\diam(S,D_h),\qquad
 a_n=\exp\{m(\log\Delta)-m(\log\Delta_n)\}.
\end{equation}
Both centers use unconditional laws, so $a_n$ is deterministic. We use
the geometric limit of \cite{MetricMaps} in the form
\begin{equation}\label{eq:metric-premise}
 a_n\longrightarrow0,
 \qquad (V_n,a_nd_n,\mu_n)
       \xrightarrow{\ d\ }(S,D_h,\mu_h)
 \quad\text{in Gromov--Hausdorff--Prokhorov topology}.
\end{equation}
The contour and topological results from that work needed in addition
to \eqref{eq:metric-premise} are stated in Section~\ref{sec:inputs}.

\subsection{Retaining the conditional path law}

Fix $T>0$. If $(K,d,\mu)$ is a compact metric probability space, let
$Q$ be a Borel probability measure on $D([0,T],K)$. We regard
$(K,d,\mu,Q)$ up to isometries transporting both measures. For two such
quadruples, define $\delta_T$ as the infimum over isometric embeddings
$i:K\to Z$ and $j:K'\to Z$ into a common complete separable metric space
of
\begin{equation}\label{eq:delta}
 \max\left\{
 d_H^Z(iK,jK'),\ d_P^Z(i_*\mu,j_*\mu'),\
 d_P^{D([0,T],Z)}(i_*Q,j_*Q')\right\}.
\end{equation}
Here $d_H$ and $d_P$ denote Hausdorff and Prokhorov distances. The last
term uses the Skorokhod $J_1$ topology: time changes are increasing
continuous bijections of $[0,T]$, and both their uniform time displacement
and the uniform spatial discrepancy tend to zero. Each of the three
distances may be truncated at one. An embedding acts on a path at every
time.

\begin{samepage}
\begin{theorem}[Linear clock and conditional laws]\label{thm:main}
For every finite $T>0$, define
\[
 Q_n^T=\Law\bigl((X^n_{nt})_{0\le t\le T}\mid M_n,T_n\bigr),
 \qquad
 Q_h^T=\Law\bigl((B^h_t)_{0\le t\le T}\mid h\bigr).
\]
With the spatial normalization \eqref{eq:spatial-normalization},
\[
 (V_n,a_nd_n,\mu_n,Q_n^T)
       \xrightarrow{\ d\ }(S,D_h,\mu_h,Q_h^T)
       \qquad\text{for }\delta_T,
\]
through all positive integers $n$. Thus the deterministic clock
multiplier for \eqref{eq:walk-generator} and \eqref{eq:continuum-form}
is $c=1$.
\end{theorem}
\end{samepage}

The random probability law is part of the state in Theorem~\ref{thm:main}.
For example, the same argument gives joint convergence of any fixed
number of conditionally independent stationary walks on the same map:
in a common embedding, the associated conditional joint laws are the
finite products of $Q_n^T$. Their weak convergence follows from that of
$Q_n^T$. This preserves information about the shared environment.

\subsection{Context and related work}
Mullin's enumeration of tree-rooted maps \cite{Mullin} gives the
combinatorial starting point: the selected spanning tree and its
complementary dual tree encode a map by a two-coordinate walk. The
continuum counterpart is the mating-of-trees construction of
Duplantier, Miller, and Sheffield \cite{DMS}. At $\gamma=\sqrt2$ its
boundary-length coordinates are independent Brownian motions. Miller
and Sheffield \cite{MSSphere} give the corresponding finite-area
sphere description, and their imaginary-geometry construction
\cite{IG4} supplies the space-filling $\mathrm{SLE}_8$ between the two
continuum trees. Couplings of discrete and continuum encodings have
also been used to compare graph distances across random-map models
\cite{GHS}. The continuum metric itself is characterized in
\cite{GMMetric}, with conformal covariance established in
\cite{GMConformal}.

Liouville Brownian motion was constructed by Garban, Rhodes, and Vargas
\cite{GRV} and independently by Berestycki \cite{Berestycki} through a
change of Brownian time determined by quantum area. Its Dirichlet-form
description \cite{GRVForm} separates two structures: the energy is
conformally invariant, while the speed measure is quantum area. This
distinction explains why the clock is a separate part of a random-map
scaling limit.

Several works establish important parts of the discrete picture.
Gwynne, Miller, and Sheffield prove convergence of the Tutte embedding
of mated-CRT maps, including conditional walk-law convergence modulo
time parameterization \cite{GMS}. Their estimates for harmonic
functions compare discrete and continuum energies up to constants and control
oscillation \cite{GMSHarmonic}; their invariance principle for ergodic
scale-free environments develops a broader framework for embedded
networks \cite{GMS2}. Gwynne and Miller obtain resistance,
return-probability, and displacement estimates for several infinite
random planar maps, including spanning-tree-weighted maps \cite{GM}.

Berestycki and Gwynne prove parameterized, quenched convergence to
Liouville Brownian motion for mated-CRT maps \cite{BG}. Their
deterministic time scale is an annealed exit-time median, comparable
to the inverse quantum area of a cell. Whether its ratio to that
inverse area converges is left open in \cite[Remark~1.3]{BG}.
Theorem~\ref{thm:main} concerns
the finite combinatorial tree-weighted map and its stationary
half-edge walk. It retains the conditional path law together with
metric and area, and identifies the acceleration as exactly $n$
for \eqref{eq:continuum-form}.

Our geometric inputs are the companion results in \cite{MetricMaps}.
The nondegeneration argument adapts the local-flow method in
\cite[Section~``The canonical conformal structure'']{ConformalMaps}
to a primal and a dual electrical network. The use of planar duality
to determine an effective conductivity has classical precedents,
including Dykhne's exact-conductivity argument \cite{Dykhne}. Here
we must first construct the energy limit without regularity
assumptions on the shapes of discrete faces. A separate Green
estimate in the original finite-map law then controls sources of
small stationary mass and supplies the missing clock information.


\subsection{Proof strategy and reusable estimates}

For a finite graph embedded in the sphere, write
\begin{equation}\label{eq:network-energy}
 \en(u)=\sum_{ab\in E}(u(a)-u(b))^2,
\end{equation}
with one term for each unoriented edge, including multiplicities. Loops
contribute zero. The same convention defines $\enstar$ on the dual graph.
The two main tasks are to identify the limit of these unnormalized
energies and to pass from that limit to motion in actual time.

The contour description of the map gives local finite graph pieces in
conformal coordinates. Section~\ref{sec:inputs} explains why the local
reconstruction and conditional-kernel calculations of \cite{MetricMaps}
also apply to extremal lengths and constrained energy minima on these
pieces. Exact edge incidences matter here: approximation of graph
distances alone would not identify an energy. The extension retains
compact numerical outputs before extracting a subsequence, and uses a
change of measure depending only on the contour inputs.

Section~\ref{sec:nondegeneration} proves bounded-energy cutoff estimates
without geometric control of individual faces. Its central argument is
a planar obstruction to simultaneous degeneracy of the primal and dual
networks. If annular extremal lengths could vanish, low-energy flows
would have a nonzero limiting average length measure. Locality and
triviality of field point germs produce transverse primal and dual
crossing flows at a typical point. Every such pair crosses an edge and
its dual, contradicting the vanishing traffic norms. The cutoffs then
yield harmonic compactness.

Section~\ref{sec:energy} identifies every local variational limit as
$\lambda\int|\nabla u|^2$, with the same positive deterministic $\lambda$
for the primal and dual. To determine $\lambda$, we compare a condenser
on an annulus with the period of its rotated current. The comparison in
one direction gives $\lambda\le1$; recovery of an angular cochain and
the signed primal--dual pairing give $\lambda\ge1$. Thus the coefficient
is exactly one in every subsequence. Section~\ref{sec:sphere} transfers
the result to fixed-area spheres using interior excursion densities and
removes the single omitted contour endpoint by zero Sobolev capacity.

The remaining speed-measure issue is addressed by the following estimate,
proved in Section~\ref{sec:green}. Let $H_n$ be the positive unweighted
graph Laplacian. For $\mu_n(f)=0$, define $G_nf$ by
$H_nG_nf=\mu_nf$ and $\mu_n(G_nf)=0$, where the product $\mu_nf$ on the
right denotes its vector of vertex masses. For every $0<\rho<1/2$ there
are map-measurable random variables $A_n$ with
$\sup_n\E A_n<\infty$ such that, simultaneously for $\|f\|_\infty\le1$,
\begin{equation}\label{eq:intro-green}
 \|G_nf\|_\infty\le A_n\,\mu_n(|f|)^\rho.
\end{equation}
Its proof is entirely in the fixed-size map law. A spanning-forest
identity expresses the centered inverse Laplacian through tree cuts;
the contour bijection expresses each cut size as one plus a dual-tree
distance. Dyck-excursion estimates then control the interval sums.
This estimate can be used independently of the energy identification.

Finally, Section~\ref{sec:dynamics} uses \eqref{eq:intro-green} to obtain
uniform resolvent convergence and a stopping-time tightness estimate.
The resolvents identify the conditional path laws by their integrated
finite-dimensional distributions. The normalization is visible already
at the discrete level:
\[
 -\langle f,nL_nf\rangle_{L^2(\mu_n)}=\tfrac12\en(f).
\]
Combining this identity with the coefficient $\lambda=1$ gives the clock
in Theorem~\ref{thm:main}.
